\chapter{Physical Therapy}

\section*{Overview}
Physical therapy is a rehabilitation service that improves movement, strength, balance, endurance, and function after injury, illness, surgery, or disability.

\section*{Why It Is Done}
It is commonly used for arthritis, back and neck pain, fractures, sports injuries, stroke, balance problems, and recovery after illness or surgery. Treatment is tailored to the person's goals and physical limitations.

\section*{How It Is Performed}
The therapist evaluates movement, strength, flexibility, balance, and functional tasks. Treatment may include therapeutic exercise, stretching, manual therapy, gait training, balance training, education, and selected modalities such as heat or cold.

\section*{Preparation and Safety}
Wear comfortable clothing and bring information about symptoms, medicines, and activities that are difficult. Exercises should be performed as instructed and progressed gradually. Report chest pain, severe shortness of breath, new weakness, or worsening pain.

\section*{Results and Follow-up}
Progress is measured by improvements in mobility, strength, pain, balance, and daily function. The therapist may prescribe a home exercise program and coordinate care with other clinicians.

\section*{When to Seek Medical Care}
Follow the procedure team's specific instructions. Seek urgent medical assessment for severe or worsening pain, uncontrolled bleeding, fever or signs of infection, trouble breathing, chest pain, fainting, new weakness or confusion, inability to urinate, or any rapidly worsening symptom. The appropriate warning signs depend on the procedure and the patient's medical condition.

\section*{References}
\begin{enumerate}
  \item MedlinePlus. Physical therapy. U.S. National Library of Medicine; 2025.
  \item Current Medical Diagnosis and Treatment. McGraw-Hill; 2025.
\end{enumerate}
\clearpage
